%% kapa: exam=1 round=27 year=2016 period=2 form=A subject=appraisal_law subject_ko=감정평가관계법규 seq=9 number=9
%% source: Q-Net 감정평가사 기출문제 (cst00309) artlSeq=5169402 file=extracted/1008/5169402/제27회 감정평가사 1차 문제/제27회 감정평가사 2교시문제(A형).pdf
%% conversion: kordoc 4.24.51 -> tools/gen1.py (auto); review=proofread-sonnet; flags=none
\begin{question}{9}
국토의 계획 및 이용에 관한 법령상 용도지역 미세분 지역에 관한 설명이다. 다음 ( )안의 내용이 옳게 연결된 것은?
\begin{framedtext}
도시지역 또는 관리지역이 세부 용도지역으로 지정되지 아니한 경우, 용도 지역별 건축물의 건축제한에 관한 규정을 적용할 때에 해당 용도지역이 도시지역인 경우에는 ( ㄱ )에 관한 규정을 적용하고, 관리지역인 경우에는 ( ㄴ )에 관한 규정을 적용한다. \par
\end{framedtext}
\begin{choices}
  \item ㄱ: 생산녹지지역, ㄴ: 보전관리지역
  \item ㄱ: 생산녹지지역, ㄴ: 자연환경보전지역
  \item ㄱ: 보전녹지지역, ㄴ: 보전관리지역
  \item ㄱ: 보전녹지지역, ㄴ: 자연환경보전지역
  \item ㄱ: 자연환경보전지역, ㄴ: 자연환경보전지역
\end{choices}
\end{question}
\answer{3}
